\section{Programma}

Di seguito viene riportato il codice sorgente del programma utilizzato per la trattazione.


\subsection{Contratto Parameters}
Per il corretto utilizzo del programma è necessario creare una classe che estende \textbf{AbstractParam}:
\begin{lstlisting}[style=myStile, caption={Parameters}, label={script:1}]
class AbstractParam(ABC):
    """
    Contratto obbligatorio per i Parametri
    """

    # ===== PARAMETRI BASE =====
    
    @property
    @abstractmethod
    def n_roads(self) -> int:
        pass

    @property
    @abstractmethod
    def n_routes_hat(self) -> int:
        pass

    @property
    @abstractmethod
    def n_routes_check(self) -> int:
        pass

    @property
    @abstractmethod
    def operation(self) -> Operation:
        pass

    @property
    def show_result(self) -> bool:
        return False
    
    @property
    def show_iterations(self) -> bool:
        return True

    #Utilizzato solo se self.operation==Operation.NASH_EQ
    @property
    def print_as_fraction(self) -> bool:
        return False

    # ===== MATRICI =====
    
    @property
    @abstractmethod
    def Gamma_hat(self) -> SafeArray:
        pass

    @property
    @abstractmethod
    def Gamma_check(self) -> SafeArray:
        pass

    @property
    @abstractmethod
    def variation_hat(self) -> SafeArray:
        pass

    @property
    @abstractmethod
    def variation_check(self) -> SafeArray:
        pass

    # ===== COSTI =====
    
    @abstractmethod
    def tau_hat(self, eta_hat, eta_check) -> SafeArray:
        """
        Ritorna il costo delle strade per la popolazione hat
        """
        pass

    @abstractmethod
    def tau_check(self, eta_hat, eta_check) -> SafeArray:
        """
        Ritorna il costo delle strade per la popolazione check
        """
        pass
    
    def tau_hat_variated(self, eta_hat, eta_check, variation) -> SafeArray:
        return self.tau_hat(eta_hat, eta_check) + self.variation_hat * variation

    def tau_check_variated(self, eta_hat, eta_check, variation) -> SafeArray:
        return self.tau_check(eta_hat, eta_check) + self.variation_check * variation
\end{lstlisting}
Le proprietà e i metodi contrassegnati da \textbf{@abstractmethod} devono essere istanziati nella classe dei parametri creata dall'utente, mentre gli altri
hanno valori base che possono essere reistanziati o meno.

\subsection{Utility}
Viene introdotta una Enum per permettere all'utente di definire quale operazione eseguire (tramite il Listato~\ref{script:1}): calcolo dell'equilibrio di Nash oppure computazione del grafico
dei tempi di percorrenza della rete dipendenti dalla variazione dei costi delle strade definiti nel Listato~\ref{script:1}.

\begin{lstlisting}[style=myStile, caption={Operation}, label={script:2}]
  class Operation(Enum):
    NASH_EQ = 1 #Calcolo equilibrio di Nash
    NASH_EQ_VARIATIONS = 2 #Computazione del Grafico delle variazioni
\end{lstlisting}

Estendiamo anche la classe \textbf{np.ndarray} (vedi \cite{numpy_docs}) per gestire i casi in cui un valore di un array/matrice posto ad \textbf{np.inf} viene moltiplicato per 0.
Ciò può capitare soprattutto durante il calcolo dell'Equazione \ref{eq:2}

\begin{lstlisting}[style=myStile, caption={Operation}, label={script:3}]
  class SafeArray(np.ndarray):
    def __new__(cls, input_array):
        obj = np.asarray(input_array, dtype=float).view(cls)
        return obj
    
    def __matmul__(self, other):
        selfChanged = False
        otherChanged = True        
        
        if np.isinf(self).any():
            self[np.isinf(self)] = 1.e200
            selfChanged = True
        if np.isinf(other).any():
            other[np.isinf(other)] = 1.e200
            otherChanged = False
            
        result = np.matmul(self, other)
        
        result[result > 1.e100] = np.inf
        
        if selfChanged:
            self[self>1.e100] = np.inf
        if otherChanged:
            other[other>1.e100] = np.inf

        return result.view(type(self))
\end{lstlisting}


\subsection{Classe Initializer}
Una volta creata una classe che estende \textbf{AbstractParam}, per eseguire il programma è sufficiente creare un'istanza di \textbf{Initializer} e invocare il metodo 
\textbf{Initializer::start}:

\begin{lstlisting}[style=myStile, caption={Initializer}, label={script:4}]
class Initializer:
    
    def __init__(self, param: AbstractParam):
        np.seterr(divide='ignore')
        
        if not isinstance(param, AbstractParam):
            raise TypeError("param deve essere un'istanza di BaseParameters")
        
        self.param = param
        self.computator = Computator(self.param)
        self.initial_theta_hat = SafeArray(np.ones((self.param.n_routes_hat, 1))*(1/self.param.n_routes_hat))
        self.initial_theta_check = SafeArray(np.ones((self.param.n_routes_check, 1))*(1/self.param.n_routes_check))
        self.limit_hat = 1e-11
        self.limit_check = 1e-11
        
        
    def start(self):
        """
            In base a operation in param, decide se calcolare l'equilibrio di Nash (senza variazioni) o costruire il grafico
            variando i costi delle strade
        """
        if self.param.operation == Operation.NASH_EQ:
            return self.getNashEquilibria()
        elif self.param.operation == Operation.NASH_EQ_VARIATIONS:
            return self.graphVariations()
\end{lstlisting}

Il metodo \textbf{Initializer::getNashEquilibria} calcola l'equilibrio di Nash

\subsection{Calcolo dell'equilibrio di Nash}

Per calcolare l'equlibrio di Nash si sfrutta la dimostrazione, \cite[Cap. 4]{MR4911797}, del Teorema \ref{th:1}, ovvero cercando un punto fisso della funzione $\mathcal{F}$.
Per fare ciò si parte da un $(\widehat{\theta}, \widecheck{\theta})\in S^{\widehat{n}}\times S^{\widecheck{\theta}}$ generico. Si ricava $\mathcal{F}(\widehat{\theta}, \widecheck{\theta})$,
che va a sostituire il valore di $(\widehat{\theta}, \widecheck{\theta})$, salvato invece come $(\widehat{\theta}_{prev}, \widecheck{\theta}_{prev})$. Questa operazione viene ripetuta finché
$\norm{\widehat{\theta} - \widehat{\theta}_{prev}} < \widehat{limit}$ e $\norm{\widecheck{\theta} - \widecheck{\theta}_{prev}} < \widecheck{limit}$, dove i $limit$ sono limiti di precisione
definiti in !!!!!. È utile introdurre anche un coefficiente di variazione, utile per !!!!, che verrà applicato a !!!!.

\begin{lstlisting}[style=myStile, caption={Equilibrio di Nash}, label={script:5}]
def getNashEquilibriaVariation(self, theta_hat, theta_check, limit_hat, limit_check, variation):
        prev_theta_hat = SafeArray(np.ones((self.param.n_routes_hat, 1))*np.inf)
        prev_theta_check = SafeArray(np.ones((self.param.n_routes_check, 1))*np.inf)
        
        count = 0
        
        while np.linalg.norm(theta_hat-prev_theta_hat) > limit_hat or np.linalg.norm(theta_check-prev_theta_check) > limit_check:
            prev_theta_hat[:] = theta_hat
            prev_theta_check[:] = theta_check
            
            theta_hat = self.f_hat(theta_hat, theta_check, variation)
            theta_check = self.f_check(theta_hat, theta_check, variation)
            count+=1
        
        if self.param.show_iterations: print("Iterazioni: " + str(count))
        return theta_hat, theta_check
\end{lstlisting}

